\textit{Nota:} el algoritmo Z (Z-function) ya está en la sección de Algoritmos Ad-Hoc; es útil combinarlo con estos.

\textit{¿Cuál usar?} \textbf{Trie}: cuando tienes \textbf{muchas} palabras/strings y necesitas buscar por prefijo repetidamente (un solo texto/patrón no lo justifica). \textbf{Aho-Corasick}: buscar \textbf{muchos patrones a la vez} en un solo texto (Trie + enlaces de fallo). \textbf{KMP}: buscar \textbf{un} patrón en un texto, determinista y sin colisiones. \textbf{Hashing (Rabin-Karp)}: comparar substrings entre sí en \(O(1)\) o buscar patrones cuando conviene comparar cadenas como números (ej. substrings repetidos, LCP entre dos posiciones). \textbf{Manacher}: específicamente para palíndromos.
